% !PRINTED 724
\weberpage{724}

Further, $5^{\nu'} \equiv 1 \pmod{\mu}$, and since $5^{\nu'}$ is not also congruent
with 1 with respect to the modulus $m$, it follows that:
\weq{7}{5^{\nu'} \equiv 1 + \mu \pmod{m}.}
Accordingly, in case (6), 2. we have
\[
(-1)^\alpha\,5^\beta = (-1)^\alpha\,5^{\beta_1}\,5^{\nu'} \equiv (-1)^\alpha\,5^{\beta_1} + \mu \pmod{m},
\]
and consequently the absolutely least residue of $(-1)^\alpha\,5^{\beta_1}$
with respect to the modulus $m$ is here negative. Thus in case (6), 2. the value
\[
c_{\beta_1} = -(-1)^\alpha
\]
is to be set.

From this there follows, by (2) (if we again write $\beta$ for $\beta_1$):
\weq{8}{\varphi(\Theta) = \sum_{0,\,\nu'-1}^{\beta} c_\beta\,\Theta^\beta = c_0 + c_1\Theta + c_2\Theta^2 + \cdots + c_{\nu'-1}\Theta^{\nu'-1},}
and here we have
\weq{9}{c_\beta = +1 \ \text{or} \ = -1,}
according as the absolutely least residue of $5^\beta$ is positive or
negative. Accordingly $\varphi(\Theta)$ is an \gesperrt{integer of the
field} $\Omega_\nu$.

It still remains to determine the behaviour of this number with
respect to the number 2, or rather with respect to the prime factor
$(1 - \Theta)$ of 2 contained in $\Omega_\nu$. If for this purpose we form
\weq{10}{(1 - \Theta)\,\varphi(\Theta) = 2\,\psi(\Theta),}
there results
\weq{11}{\psi(\Theta) = \frac{c_0 + c_{\nu'-1}}{2} + \frac{c_1 - c_0}{2}\,\Theta + \frac{c_2 - c_1}{2}\,\Theta^2 + \cdots + \frac{c_{\nu'-1} - c_{\nu'-2}}{2}\,\Theta^{\nu'-1}.}

The coefficients in this expression are all $= 0$ or
$= 1$, and therefore $\psi(\Theta)$ is likewise an integer of the field $\Omega_\nu$.
For this number, however, if one sets $\Theta \equiv 1$, there results:
\weq{12}{\psi(\Theta) \equiv c_{\nu'-1} \equiv \pm 1\ [\mathrm{mod}\ (1-\Theta)],}
and consequently $\psi(\Theta)$ is relatively prime to 2.

The product occurring in (4) is nothing other than the norm of the
number $\varphi(\Theta)$ taken with respect to the field $\Omega_\nu$,
and if we denote this norm by $N$, then by \S.\,169
\[
N(1-\Theta) = 2,
\]
%CONT>>
